\documentclass{report}
\usepackage{amsmath,amssymb}
\setlength{\parindent}{0mm}
\setlength{\parskip}{1em}
\begin{document}
\begin{center}
\rule{6in}{1pt} \
{\large Christopher Siefert \\
{\bf Mimetic Least Squares Methods with Preconditioners for Darcy Flow}}

Sandia National Laboratories \\ P O Box 5800 \\ MS 1323 \\ Albuquerque \\ NM 87185-1323
\\
{\tt csiefer@sandia.gov}\\
Pavel Bochev\\
Oksana Guba\end{center}

We present a mimetic least squares method for Darcy flow,
using nodal elements for pressure and lowest order face elements for the
fluxes. Because the pressure and flux are not subject to a joint inf-sup
condition, we can choose the discrete spaces independently. The choice of
a face element discretization for the fluxes gives us a method with
improved stability, accuracy and conservation properties.

The mimetic least-squares functional is norm-equivalent to a norm on
$H^1(\Omega)\times H(div)$. As a result, the associated linear system,
though positive definite, has a 2x2 block structure, where one diagonal
element is a nodal Laplacian, and the other is a grad-div problem. Owing
to the norm-equivalence of the least-squares functional, the system can
be effectively preconditioned by its diagonal.

We treat the least-squares algebraic equations with smoothed
aggregation algebraic multigrid and the latter with the compatible
gauge approach of Bochev et al. [1]. The structure of the equations
allows us to combine these into a simple block preconditioner. We
demonstrate good algorithmic and parallel scalability of our problem on
up to 17,000 cores on NERSC's hopper machine.

[1] P. Bochev, C. Siefert, R. Tuminaro, J. Xu and Y. Zhu. Compatible
Gauge Approaches for H(div) Equations. Technical Report, SAND 2007-5384P,
Sandia National Laboratories, August 2007.


\end{document}
